\section{Reproducibility}
\label{app:repro}

\subsection{Fixed source}\label{fixed-source}

\begin{itemize}
\tightlist
\item
  PrimeKG, Harvard Dataverse DOI: 10.7910/DVN/IXA7BM.
\item
  File ID: 6180620; filename: kg.csv.
\item
  Download: https://dataverse.harvard.edu/api/access/datafile/6180620 .
\item
  Size: 981,751,236 bytes.
\item
  Official MD5, verified: aac8191d4fbc5bf09cdf8c3c78b4e75f.
\item
  SHA-256:
  0a5a010fcd88dcec572a573d98c26c1c244a11567092b6161d3ded611cfb5c69.
\item
  Source metadata and processing audit are in data/raw; raw CSV is
  excluded from the handoff archive and version control.
\end{itemize}

\subsection{Reproduction}\label{reproduction}

Create a Python 3.11 environment and install requirements-lock.txt. In
this workspace the interpreter is \texttt{.venv/Scripts/python.exe}
(3.11.9). Run from the repository root:

\begin{verbatim}
python -m unittest discover -s tests -v
python scripts/download_data.py
python scripts/run_a_pipeline.py
\end{verbatim}

The theory, prediction and null-model results (Sections~\ref{sec:theory}
and~\ref{sec:prediction}) are regenerated from the frozen graphs and the stored
WS2, WS3 and WS5 results with:

\begin{verbatim}
python -m src.run_ws4 primekg_5000_bfs_v1 primekg_2000_bfs_v1
python -m src.run_ws4_report primekg_5000_bfs_v1
python -m src.run_null_model primekg_5000_bfs_v1
python -m src.plot_ws4 primekg_5000_bfs_v1
python -m src.plot_null_model primekg_5000_bfs_v1
\end{verbatim}

All outputs are written to \texttt{results/ws4/} and \texttt{figures/act1/} or
\texttt{figures/act2/}; random seeds are fixed in the scripts. The null-model
run takes a few minutes on a laptop.

The downloader requires data/raw/dataverse\_metadata.json; README gives
the exact metadata download command. It supports resume and verifies
each range response plus the official checksum. Pipeline stages refuse
to overwrite frozen core graph files. To rebuild without existing
outputs, use a fresh directory with the code and raw source, or a
separate G17\_ROOT. \texttt{scripts/check\_reproducibility.py}
demonstrates an isolated build and compares scientific-file hashes.

The final formal runs use Python 3.11.9 with requirements-lock.txt.
Initial graph generation used Python 3.14.3; manifests preserve this
history. Rebuilding under 3.11.9 reproduced nodes, edges, queries and
random orders exactly: 68 scientific files match SHA-256. Environment
fields and other descriptive metadata are not claimed byte-identical
across interpreters. Source fingerprints in
results/ws1/*\_\_provenance.json connect each result to its data and
algorithm.

\subsection{Determinism}\label{determinism}

Original node IDs are parsed as integers. BFS starts at original disease
ID 32617 and visits neighbors in ascending ID order. The subgraph is
induced. Nodes are assigned local IDs in ascending original-ID order.
Edges are sorted by relation,u,v before assigning eid. Queries are
sorted before independent seed-42 group permutations; selected index
positions are re-sorted. The training split uses a separate
default\_rng(42) permutation. Damage seeds are 0--9 and masks use fixed
permutation prefixes with round(f*total).

Cross-graph entity joins use x\_index. Random orders belong to one gid
and one ID universe. Library versions matter for sampled centrality; B
should record them separately. A's frozen manifest excludes future
B-owned centrality files and should not be rewritten to register them.

\subsection{Validation evidence}\label{validation-evidence}

\begin{itemize}
\tightlist
\item
  results/validation/reproducibility.json: 68 byte-identical files and
  BFS prefix check.
\item
  tests/test\_percolation.py: all-state comparison on 80 small graphs,
  semantic duplicates, edge cases and input validation.
\item
  tests/test\_theory.py: exact 3-regular formula and isolated/empty
  degree cases.
\item
  results/validation/\{gid\}\_\_networkx\_checks.csv: 10 real-graph
  states per gid, six counts per state.
\item
  results/validation/\{gid\}\_\_retrieval\_reference.csv: independent
  semantic-path/distance reference for downstream consumers; not a claim
  that C's module has been tested.
\item
  tests/test\_theory\_d.py and tests/test\_null\_model.py: analytic cases of
  the retrieval and GF formulas; degree, typed-degree and joint-degree
  preservation of the three rewirings.
\item
  results/ws4/\{gid\}\_\_configuration\_model\_check.csv,
  \_\_monte\_carlo\_check.csv, \_\_uncertainty.csv,
  \_\_mincut\_Linf.csv and \_\_restoration\_prediction.csv: the
  agreement figures quoted in Sections~\ref{sec:theory} and~\ref{sec:prediction}.
\item
  slides/A\_slides\_3\_4.pptx: two editable slides with native charts
  and an editable table; the scientific PNG/PDF plots retain shaded SD
  bands, while the native slides show SD bounds.
\end{itemize}

\subsection{Limits and
responsibilities}\label{limits-and-responsibilities}

The full graph's queries.json is intentionally empty because this
release defines only structural full-graph comparisons. It is not an R=0
dataset. The 2,000-node graph contains 112 queries and only three drugs.
Report functions must handle actual group sizes. Query splits allow
entity overlap.

A's isolated GF reference and D's final module \texttt{src/theory.py} agree
to within $2\times10^{-13}$ on a 201-point grid for both bond and site
percolation. E
can supply plotting.py without changing scientific data. Full group
integration, rehearsal, institutional submission and peer grading are
not performed by these scripts.
